\documentclass[11pt,a4paper]{article}
\usepackage[margin=1.25cm, top=1.1cm, bottom=1.2cm]{geometry}
\usepackage[scaled]{helvet}\renewcommand\familydefault{\sfdefault}
\usepackage[T1]{fontenc}
\usepackage{tikz,xcolor,array,tabularx,booktabs,enumitem,multicol}
\usepackage{amsmath,amssymb}
\usetikzlibrary{arrows.meta,shapes.misc,calc}
\pagestyle{empty}\setlength{\parindent}{0pt}
\definecolor{ttdark}{HTML}{2B2D5C}
\definecolor{ttaccent}{HTML}{D9480F}
\definecolor{ttglow}{HTML}{25B5A4}
\definecolor{ttpale}{HTML}{EEF1FA}
\tikzset{flag/.style={draw=ttaccent, fill=white, text=ttaccent, line width=0.8pt, font=\tiny\bfseries, rounded corners=2pt, inner sep=2pt, align=center}}
\newcommand{\banner}[1]{{\color{ttdark}\large\bfseries #1}\par\vspace{-7pt}{\color{ttdark}\rule{\linewidth}{1.2pt}}\par\vspace{-2pt}}
\newcommand{\dial}[3]{\tikz[baseline=-3.5pt]{\node[draw=ttdark, line width=1pt, rounded corners=3pt, inner sep=1.5pt] {%
  \tikz[baseline]{\node[draw=black!45, text=ttdark, font=\scriptsize\bfseries, inner sep=2.5pt, minimum width=1.5cm]{#1};}\,%
  \tikz[baseline]{\node[draw=black!45, text=black, font=\scriptsize\bfseries, inner sep=2.5pt, minimum width=1.3cm]{#2};}\,%
  \tikz[baseline]{\node[draw=black!45, text=ttdark, font=\scriptsize\bfseries, inner sep=2.5pt, minimum width=1.6cm]{#3};}};}}
\newcommand{\aline}{\rule{2.6cm}{0.4pt}}
\newcommand{\lline}{\rule{4.6cm}{0.4pt}}
\newcommand{\ans}[1]{\textcolor{red!75!black}{\bfseries #1}}
\newcommand{\factbox}[2]{\begin{tikzpicture}\node[draw=ttdark, rounded corners=4pt, line width=1pt, inner sep=7pt, text width=#1, align=left, font=\small]{#2};\end{tikzpicture}}
\begin{document}

\begin{center}
{\LARGE\bfseries Worked Solutions}\\[2pt]
{\large Time Travellers! --- Calendars, conversions and timelines}
\end{center}
\vspace{2pt}

\banner{Part A \quad Calendar hops --- small jumps}
Starting point: \textbf{Saturday 31 October 2026, 10:00}.
\begin{enumerate}[label=\textbf{A\arabic*.}, leftmargin=*, itemsep=6pt, topsep=4pt]
\item \dial{FORWARD}{1}{WEEK}. \(1\) week \(=7\) days. \(31\text{ Oct}+7\text{ days}=7\text{ Nov}\). The day is still Saturday. \ans{Saturday 7 November 2026}.

\item \dial{FORWARD}{10}{DAYS}. \(10=7+3\), so the day moves forward \(3\) days from Saturday to Tuesday. The date is \(31\text{ Oct}+10\text{ days}=10\text{ Nov}\). \ans{Tuesday 10 November 2026}.

\item \dial{BACK}{2}{WEEKS}. \(2\) weeks \(=14\) days. \(31\text{ Oct}-14\text{ days}=17\text{ Oct}\). The day is still Saturday. \ans{Saturday 17 October 2026}.

\item \dial{FORWARD}{3 WEEKS}{+ 2 DAYS}. \(3\) weeks \(=21\) days, plus \(2\) days \(=23\) days. \(31\text{ Oct}+23\text{ days}=23\text{ Nov}\). Day: Saturday \(+23\) days; \(23\equiv 2\pmod 7\), so Monday. \ans{Monday 23 November 2026}.

\item \dial{BACK}{45}{DAYS}. \(45=31+14\). \(31\text{ Oct}-31\text{ days}=30\text{ Sep}\). Then \(30\text{ Sep}-14\text{ days}=16\text{ Sep}\). Day: Saturday \(-45\) days; \(45\equiv 3\pmod 7\), so Wednesday. \ans{Wednesday 16 September 2026}.

\item \dial{FORWARD}{100}{HOURS}. \(100\text{ h}=4\text{ days }4\text{ h}\). \(31\text{ Oct }10{:}00+4\text{ days}=4\text{ Nov }10{:}00\). Add \(4\text{ h}\), giving \(14{:}00\). Day: Saturday \(+4\) days \(=\) Wednesday. \ans{Wednesday 4 November 2026 at 14:00}.

\item Bonfire Night (5 November) to Christmas Day (25 December). November: \(5\to 30\) is \(25\) days. December: \(1\to 25\) is \(25\) days. Total \(25+25=50\) days. Since \(50=7\times 7+1\), this is \(7\) weeks \(1\) day. \ans{50 days = 7 weeks 1 day}.

\item Whole weeks add multiples of \(7\) days. From Saturday 31 October, every whole-week jump lands on a Saturday. The gap to Christmas Day is
\[
31\text{ Oct}\to 30\text{ Nov}=30\text{ days},\qquad
30\text{ Nov}\to 25\text{ Dec}=25\text{ days},
\]
so the total is \(55\) days. This is not a multiple of \(7\). The nearest whole-week landings are \(7\) weeks \(=49\) days later, Saturday 19 December, and \(8\) weeks \(=56\) days later, Saturday 26 December. \ans{No. Nearest reachable Saturdays: 19 December 2026 and 26 December 2026}.

\item Check the calendars for a Friday the 13th. The only one in September--December 2026 is 13 November. \ans{Friday 13 November 2026}.
\end{enumerate}

\banner{Part B \quad Conversion station --- medium jumps}
\textbf{1. Conversions.}
\begin{itemize}[leftmargin=*, itemsep=2pt, topsep=2pt]
\item \(3\text{ weeks}=3\times 7=21\text{ days}\).
\item \(4\text{ centuries}=4\times 10=40\text{ decades}\).
\item \(2\text{ days}=2\times 24=48\text{ hours}\).
\item \(2\text{ millennia}=2\times 10=20\text{ centuries}\).
\item \(180\text{ minutes}=180\div 60=3\text{ hours}\).
\item \(7200\text{ seconds}=7200\div 3600=2\text{ hours}\).
\item \(5\text{ decades}=5\times 10=50\text{ years}\).
\item \(3\text{ days}=3\times 24\times 60=4320\text{ minutes}\).
\end{itemize}

\textbf{2. Dial readings.}
\begin{enumerate}[label=\textbf{B\arabic*.}, leftmargin=*, itemsep=4pt, topsep=4pt]
\item \dial{FORWARD}{72}{HOURS}: \(72\div 24=3\) days. \ans{3 days}.
\item \dial{FORWARD}{420}{DAYS}: \(420\div 7=60\) weeks. \ans{60 weeks}.
\item \dial{BACK}{1000}{HOURS}: \(1000\div 24=41\) remainder \(16\), so \(41\) days \(16\) hours. \ans{41 days 16 hours}.
\item \dial{BACK}{1000000}{SECONDS}:
\[
1\,000\,000\div 60=16666\text{ min}\quad\text{remainder }40\text{ s},
\]
\[
16666\div 60=277\text{ h}\quad\text{remainder }46\text{ min},
\]
\[
277\div 24=11\text{ days}\quad\text{remainder }13\text{ h}.
\]
So the reading is \(11\) days \(13\) hours \(46\) minutes \(40\) seconds. \ans{11 d 13 h 46 min 40 s}.
\end{enumerate}

\textbf{3. Leap years.} A year is leap if it is divisible by \(4\), except century years, which must be divisible by \(400\).
\[
2028\div 4=507\quad\text{(leap)},\qquad
2030\nmid 4\quad\text{(not leap)},
\]
\[
2100\div 400\notin\mathbb{Z}\quad\text{(not leap)},\qquad
2400\div 400=6\quad\text{(leap)}.
\]
\ans{2028 and 2400}.

\textbf{4. Order from shortest to longest.} Convert to days:
\[
1\,000\,000\text{ s}=11\text{ d }13\text{ h }46\text{ min }40\text{ s}\approx 11.57\text{ d},
\]
\[
50\,000\text{ min}=50\,000\div 60\div 24\approx 34.72\text{ d},
\]
\[
0.1\text{ years}=0.1\times 365=36.5\text{ d},
\]
\[
1000\text{ h}=1000\div 24\approx 41.67\text{ d},
\]
\[
6\text{ weeks}=42\text{ d}.
\]
\ans{\(1\,000\,000\text{ s} < 50\,000\text{ min} < 0.1\text{ years} < 1000\text{ hours} < 6\text{ weeks}\).}

\banner{Part C \quad Timeline jumps --- giant leaps}
Recall that there is no year \(0\). For a BC year \(B\) and an AD year \(A\), the time between them is \(A+B-1\) years.

\textbf{1. Marking the big timeline.}
\begin{itemize}[leftmargin=*, itemsep=2pt, topsep=2pt]
\item Tutankhamun dies \(1323\) BC: between \(1500\) BC and \(1000\) BC, closer to \(1500\) BC.
\item First Olympic Games \(776\) BC: between \(1000\) BC and \(500\) BC.
\item Romans invade Britain AD \(43\): just to the right of the \(1\) BC/AD \(1\) mark.
\item Battle of Hastings AD \(1066\): just after AD \(1000\).
\end{itemize}
\ans{See crosses on the big timeline at these positions.}

\textbf{2. Years ago from 2026.}
\begin{enumerate}[label=(\alph*), leftmargin=*, itemsep=2pt, topsep=2pt]
\item Battle of Hastings AD \(1066\): \(2026-1066=960\) years.
\item Roman invasion AD \(43\): \(2026-43=1983\) years.
\item First Olympic Games \(776\) BC: \(2026+776-1=2801\) years.
\item Great Pyramid \(2560\) BC: \(2026+2560-1=4585\) years.
\end{enumerate}
\ans{(a) 960 \quad (b) 1983 \quad (c) 2801 \quad (d) 4585}.

\textbf{3. 2(a) in centuries and decades.} \(960=9\times 100+6\times 10\). \ans{9 centuries and 6 decades}.

\textbf{4. Landing places.}
\begin{itemize}[leftmargin=*, itemsep=2pt, topsep=2pt]
\item \dial{BACK}{3}{MILLENNIA}: \(3\) millennia \(=3000\) years. \(3000-2026+1=975\) BC. \ans{975 BC}.
\item \dial{FORWARD}{5}{CENTURIES}: \(5\) centuries \(=500\) years. \(2026+500=2526\) AD. \ans{AD 2526}.
\item \dial{BACK}{25}{CENTURIES}: \(25\) centuries \(=2500\) years. \(2500-2026+1=475\) BC. \ans{475 BC}.
\end{itemize}

\textbf{5. Zoomed timeline marks.} Place crosses at Magna Carta \(1215\), Great Fire \(1666\), Moon landing \(1969\), and Today \(2026\). \ans{See crosses on the zoomed timeline at these years.}

\textbf{6. Differences.}
\[
\text{(a) }1666-1215=451\text{ years},\qquad
\text{(b) }2026-1969=57\text{ years}.
\]
\ans{(a) 451 years \quad (b) 57 years}.

\textbf{7. Is the Great Fire closer to Hastings or to today?}
\[
\text{Great Fire }1666\text{ to Hastings }1066:\quad 1666-1066=600\text{ years},
\]
\[
\text{Great Fire }1666\text{ to today }2026:\quad 2026-1666=360\text{ years}.
\]
Since \(360<600\), the Great Fire is closer to today. \ans{Closer to today (360 years vs 600 years).}

\banner{Part D \quad Mission log --- getting home}
Battery rule: every century of travel uses \(2\%\). Equivalently,
\[
\text{battery used}=\frac{\text{years}}{100}\times 2\%=\frac{\text{years}}{50}\%.
\]

\begin{center}
\begin{tabular}{c|p{6cm}|c|c|p{3.2cm}}
\textbf{Jump} & \textbf{Dial} & \textbf{From} & \textbf{Lands in} & \textbf{Battery used}\\
\hline
1 & \dial{BACK}{3}{CENTURIES} & AD 2026 & AD 1726 & \(3\times2\%=6\%\)\\[2pt]\hline
2 & \dial{FORWARD}{25}{YEARS} & AD 1726 & AD 1751 & \(0.25\times2\%=0.5\%\)\\[2pt]\hline
3 & \dial{BACK}{2}{MILLENNIA} & AD 1751 & 250 BC & \(20\times2\%=40\%\)\\[2pt]\hline
4 & \dial{FORWARD}{1500}{YEARS} & 250 BC & AD 1251 & \(15\times2\%=30\%\)\\[2pt]\hline
\end{tabular}
\end{center}
Check: jump 1: \(2026-300=1726\). Jump 2: \(1726+25=1751\). Jump 3: \(2000-1751+1=250\) BC. Jump 4: from \(250\) BC forward \(1500\) years gives \(1500-250+1=1251\) AD.

\begin{enumerate}[label=\textbf{\arabic*.}, leftmargin=*, itemsep=6pt, topsep=4pt]
\item From AD \(1251\) to AD \(2026\): \(2026-1251=775\) years. \ans{FORWARD 775 years}.

\item Total travel:
\[
300+25+2000+1500=3825\text{ years}.
\]
In centuries: \(3825\div 100=38.25\). \ans{3825 years = 38.25 centuries}.

\item Battery used: \(38.25\times 2\%=76.5\%\). Battery left: \(100\%-76.5\%=23.5\%\). \ans{76.5\% used; 23.5\% left}.

\item Home jump \(=775\) years \(=7.75\) centuries. Battery needed:
\[
7.75\times 2\%=15.5\%.
\]
Since \(23.5\%>15.5\%\), she has enough. She arrives home with \(23.5\%-15.5\%=8\%\) left. \ans{Yes; 15.5\% needed, 8\% left on arrival}.

\item One possible mission plan:
\begin{center}
\begin{tabular}{c|p{4.5cm}|p{2.5cm}|p{3.5cm}}
\textbf{Jump} & \textbf{Dial setting} & \textbf{Lands in} & \textbf{Battery used}\\
\hline
1 & BACK \(4585\) years & 2560 BC & \(45.85\times2\%=91.7\%\)\\
2 & FORWARD \(3625\) years & AD 1066 & \(36.25\times2\%=72.5\%\)\\
3 & FORWARD \(903\) years & AD 1969 & \(9.03\times2\%=18.06\%\)\\
Home & FORWARD \(57\) years & AD 2026 & \(0.57\times2\%=1.14\%\)\\
\hline
\multicolumn{3}{r|}{\textbf{Total}} & \(183.4\%\)\\
\end{tabular}
\end{center}
Workings: to \(2560\) BC: \(2026+2560-1=4585\) years back. To AD \(1066\): \(2560+1066-1=3625\) years forward. To AD \(1969\): \(1969-1066=903\) years forward. Home: \(2026-1969=57\) years forward. Total years:
\[
4585+3625+903+57=9170=91.7\text{ centuries}.
\]
Battery needed:
\[
91.7\times2\%=183.4\%.
\]
This is more than \(100\%\), so she would need to recharge. In fact, any order visiting all three and returning needs at least \(9170\) years of travel. \ans{Plan above; total battery needed \(183.4\%\), so a recharge is required.}
\end{enumerate}

\end{document}